\chapter{Componentについて }
\hypertarget{section_component}{}\label{section_component}\index{Componentについて@{Componentについて}}
\begin{DoxyItemize}
\item 単純な機能を持ったパーツです \item Objectのつけて利用できます \item 基本同じタイプのコンポーネントは付けれません。(「\+Component\+Transform」を2つとかは付けれません)\end{DoxyItemize}
\begin{DoxyItemize}
\item コンポーネントの利用方法 
\begin{DoxyCode}{0}
\DoxyCodeLine{　}
\DoxyCodeLine{\ \ \ obj-\/>AddComponent<\mbox{\hyperlink{class_component_model}{ComponentModel}}>();\ \ \textcolor{comment}{//<\ 「ComponentModel」を\ objにて利用する}}
\DoxyCodeLine{}
\DoxyCodeLine{\ \ \ \textcolor{comment}{//\ ※objはObjectタイプで、先に次のように作成しておく必要があります【例】}}
\DoxyCodeLine{\ \ \ \textcolor{keyword}{auto}\ obj\ =\ \mbox{\hyperlink{class_scene_a8eafc112b91ebde95fd4491d409b5334}{Scene::CreateObjectPtr<Object>}}();}
\DoxyCodeLine{　}

\end{DoxyCode}
\end{DoxyItemize}
\begin{DoxyItemize}
\item 使用しなくなったコンポーネントの削除 
\begin{DoxyCode}{0}
\DoxyCodeLine{　}
\DoxyCodeLine{\ \ \ obj-\/>RemoveComponent<\mbox{\hyperlink{class_component_model}{ComponentModel}}>();\ \ \textcolor{comment}{//<\ 「ComponentModel」を\ objにて削除します}}
\DoxyCodeLine{　}

\end{DoxyCode}
 \end{DoxyItemize}
